

%----------------------------------------------------------------------------------------
%	PACKAGES AND OTHER DOCUMENT CONFIGURATIONS
%----------------------------------------------------------------------------------------

\documentclass[
	a4paper,
	10pt,
	unnumberedsections,
	twoside,
]{LTJournalArticle}

\addbibresource{sample.bib}

\runninghead{Concept Drift in Media Bias Classification}

\footertext{\textit{Machine Learning Methods for Data Streams}}
\usepackage{amsmath} 
\setcounter{page}{1}

%----------------------------------------------------------------------------------------
%	TITLE SECTION
%----------------------------------------------------------------------------------------

\title{Concept Drift Detection in Polarized News Streams:\\
An Online Machine Learning Approach}

\author{Piotr Jurczyk, Krzysztof Krawiec}

\date{\footnotesize\textsuperscript{Warsaw University of Technology, Faculty of Mathematics and Information Science}}

\renewcommand{\maketitlehookd}{%
\begin{abstract}
Online political discourse constantly evolves, creating potential concept drift that degrades static NLP models. This paper develops and evaluates two complementary approaches for drift detection in news streams sourced from the GDELT database spanning the 2024 U.S. Presidential Election. 

In the first part (Classification-based drift) we compare two vectorization strategies paired with online classifiers on synthetic data with controlled drifts, where ADWIN reliably detects abrupt shifts. However, on real-world political text, both approaches tend to collapse to majority-class prediction with no detectable correlation to major political events.

In the second part (Distributional drift) we implement an unsupervised distributional drift detection pipeline using daily aggregated article streams, cosine distance between consecutive window embeddings, and drift detectors. This approach avoids the requirement for accurate classification and proved ability to detect some signals synchronous with political events.
\end{abstract}
}

%----------------------------------------------------------------------------------------

\begin{document}

\maketitle

%----------------------------------------------------------------------------------------
%	ARTICLE CONTENTS
%----------------------------------------------------------------------------------------

\section{Introduction}

News media streams reflect rapidly evolving political language: terms that signal one ideology may be adopted by opposing factions within months. This creates \textit{concept drift}: a change in the joint distribution $P(X, y)$ over time that degrades static NLP models \cite{gama2014survey}. Batch-trained classifiers work on a fixed mapping from features to labels: when semantics shift around major political events, they cannot adapt without retraining.

This paper develops and evaluates two complementary approaches for drift detection in online news streams. Our contributions are:
\begin{enumerate}
    \item A GDELT data scraping pipeline for politically polarized articles spanning the 2024 U.S. Presidential Election.
    \item A synthetic benchmark generator simulating abrupt, gradual, and recurring drifts with known ground truth for controlled evaluation.
    \item \textbf{Part I}: A systematic comparison of classification-based drift detection using TF-IDF and dense LLM embeddings with online classifiers.
    \item \textbf{Part II}: An unsupervised distributional drift detection pipeline using cosine distances between consecutive embedding-based window representations and drift detectors.
    \item Empirical evidence that classification-based approach fails to detect clear concept drift in real-world political text, despite reliable detection on synthetic data, whereas unsupervised approach yielded interesting results and proved to be an useful tool in working with media streams.
\end{enumerate}

The remainder of this paper is organized as follows. Section 2 details the data acquisition, synthetic generation, and dynamic text vectorization techniques. Sections 3 and 4 outline the methodologies for the supervised classification and unsupervised distributional drift approaches. Section 5 presents the experimental results across both synthetic and real-world streams. Finally, Section 6 discusses the implications, limitations, and future directions.

%------------------------------------------------

\section{Methodology}

Our system is aa Online Machine Learning pipeline implemented primarily using the \texttt{river} Python library \cite{montiel2021river}, extended with custom components for vectorization and unsupervised drift detection. The pipeline processes one article at a time in strict temporal order, with test-then-train evaluation approach.

\subsection{Data Acquisition and Preprocessing}

\paragraph{GDELT-based real-world stream.}
We query the GDELT API v2 using political keywords filtered to left-leaning outlets (e.g., \texttt{thenation.com}, $y=0$) and right-leaning outlets (e.g., \texttt{breitbart.com}, $y=1$) from October 2024 until April 2026. Full article text is extracted using \texttt{trafilatura}. For the classification experiment, to mitigate class imbalance while preserving temporal order, we apply max-one-per-class-per-day sampling. Due to data availability constraints on certain days, the final dataset comprises 890 articles, compared to a theoretical maximum of approximately 1,100.

\paragraph{Synthetic data stream.}
We developed a custom class with two politically-coded vocabularies ($V_L$, $V_R$), each containing 50 terms plus a shared 130-word neutral vocabulary. Each synthetic article is generated as a bag of words: 50\% from the true-class vocabulary, 10\% cross-contamination, and 40\% neutral noise. Concept drift is simulated by swapping words between vocabularies. The generator supports three configurable drift types: abrupt (instantaneous swap at step $t$), gradual (distributed swap between $t_{\text{start}}$ and $t_{\text{end}}$), and recurring (vocabulary restoration).

\subsection{Dynamic Text Vectorization}

Online classification of text requires feature extraction that can handle an evolving and unbounded vocabulary without access to the full corpus upfront. We evaluate two fundamentally different approaches.

\paragraph{Incremental TF-IDF (baseline).}
The \texttt{river} library's \texttt{TFIDF} component maintains running term statistics, producing sparse vectors. TF-IDF is a lightweight bag-of-words baseline blind to word order and semantic similarity.

\paragraph{Dense LLM embeddings.}
We integrate the \texttt{all-MiniLM-L6-v2} Sentence Transformer \cite{wang2020minilm} to map text to a fixed 384-dimensional space where semantically similar texts are proximate. A custom bridge class inherits from \texttt{river.base.Transformer}, encoding each article into 384 named features on the fly, enabling the classifier to reason about semantic similarity across novel vocabularies.

\subsection{Part I: Adaptive Streaming Classifiers}
In the first part our focus was to build a binary classifier that takes one article at a time and predicts whether it comes from left or right-wing newspaper.
\paragraph{TF-IDF pipeline.}
The baseline vectorizer is composed with a \texttt{MultinomialNB} classifier, the standard choice for sparse count-based features. This pipeline was evaluated on both the synthetic stream and the real-world GDELT corpus. Drift detection used the ADWIN (Adaptive Windowing) algorithm, monitoring per-sample binary classification error.

\paragraph{LLM pipeline.}
The \texttt{StreamSentenceTransformer} is composed with \texttt{StandardScaler} (normalizing dimensions online) and \texttt{LogisticRegression}. Unlike tree-based methods with grace periods, SGD updates immediately from the first observation, essential in low-volume streams.

\paragraph{Majority-class baseline.}
For each position in the stream, we identify the most frequent label within the preceding 50 observations and compute the accuracy of predicting that class, distinguishing genuine learning from degenerate behavior.

\subsubsection{Drift Detection.}

For the synthetic experiments we implement a \textit{reset-on-drift} strategy using ADWIN: upon an alert, the active model is instantiated fresh while the previous model is retained as a \textit{shadow model} evaluated in parallel. Shadow models receive no further training or continue to train alongside the active model, allowing direct measurement of how quickly a fresh model adapts relative to its predecessor.

\subsection{Part II: Unsupervised Distributional Drift Detection}

The second part takes a complementary approach by abandoning the requirement for accurate binary classification. Instead, we detect drift by measuring semantic shifts in aggregated news streams without relying on class labels.

\subsubsection{Time-windowed stream aggregation.}

Rather than processing articles individually, articles are grouped into weekly or daily aggregates. For each time window and stream group (mixed, left-leaning, right-leaning), all article texts are concatenated to produce a single document per window. This approach reduces noise from individual articles and captures the aggregate semantic state of political media during each period.

\subsubsection{Embedding-based distance computation.}

Each aggregated window document is vectorized using the \texttt{all-MiniLM-L6-v2} Sentence Transformer, producing a 384-dimensional embedding. The cosine distance between consecutive window embeddings is computed, yielding a univariate \textit{drift score stream}: $d_i = \text{cosine\_distance}(\mathbf{e}_i, \mathbf{e}_{i+1})$ where $\mathbf{e}_i$ is the embedding for window $i$. High cosine distances indicate large shifts in the semantic composition of political discourse.

\subsubsection{Multiple drift detectors}

Three drift detectors from the \texttt{river} library are applied independently to the cosine distance stream:
\begin{itemize}
    \item \textbf{ADWIN}: Detects distributional shifts using a sliding window with $\delta = 0.1$.
    \item \textbf{PageHinkley}: Monitors cumulative differences relative to baseline with $\delta = 0.0005$, $\text{threshold} = 3$, $\alpha = 0.99$.
    \item \textbf{KSWIN}: Kernel-based statistical window test with $\alpha = 0.05$, adaptive window size.
\end{itemize}

Drift detections are matched against known political events (election, inauguration, debates, trials, etc.) within a tolerance window of 7--14 days, depending on whether daily or weekly aggregation is used.

%------------------------------------------------

\section{Conducted Experiments and Results}


\subsection{Synthetic Streams with Controlled Drift}



\paragraph{Basic drift simulation.}
A stream of 3,200 articles is generated with abrupt drifts at steps 800, 1,600, and 2,400 (swapping 50\% of vocabulary). The TF-IDF + MultinomialNB pipeline with reset-on-drift strategy (all models continue training) shows rolling accuracy climbing during stable segments, dropping sharply at each drift point, and recovering after ADWIN detection and model reset (Figure~\ref{fig:synthetic}).
\begin{figure*}
    \centering
    \includegraphics[width=\linewidth]{Report/synthetic1.png} 
    \caption{Performance on the basic synthetic stream with three abrupt concept drifts. ADWIN alerts (red dashed lines) closely follow each injected drift (purple solid lines) and trigger model resets. Fresh models (later-colored traces) recover to near-100\% accuracy within hundreds of steps, while frozen shadow models (earlier traces) show degraded performance after each subsequent drift event.}
    \label{fig:synthetic}
\end{figure*}

\paragraph{Advanced drift simulation.}
A 4,000-article stream features: (1) abrupt drift at step 800 (swap 50\%), (2) gradual drift steps 1,500--2,000 (swap 90\%), and (3) recurring concept at step 3,000 (restore original). Under \texttt{active\_only} strategy, shadow models freeze at replacement. Gradual drift does not trigger any alerts. After restoring the original vocabulary, the frozen model achieves perfect accuracy faster than newly trained one (Figure~\ref{fig:synthetic2}).

\begin{figure*}
    \centering
    \includegraphics[width=\linewidth]{Report/synthetic2.png} 
    \caption{Performance on the advanced synthetic stream with mixed drift types. The model encounters an abrupt drift at step 800, a gradual drift spanning steps 1,500--2,000, and a recurring concept drift at step 3,000 (vocabulary restoration). ADWIN detection is less responsive to gradual drift due to slow evidence accumulation, but sharply alerts on the recurring drift when pre-adapted weights fail against restored original vocabulary.}
    \label{fig:synthetic2}
\end{figure*}

\subsection{Real-World Stream Classification}
\paragraph{TF-IDF Word Representation}
The 890 GDELT articles are processed chronologically using the max-one-per-class-per-day sampling strategy. The TF-IDF + MultinomialNB pipeline exhibits rolling accuracy that precisely tracks the majority-class baseline throughout the stream (Figure~\ref{fig:tfidf}), indicating complete failure to extract any discriminative signal from political vocabulary. 

ADWIN drift detection similarly fails to trigger alerts despite two major political events: the November 5, 2024 U.S. Presidential Election and the January 20, 2025 Inauguration. The overlap between model accuracy and majority-class baseline reveals that TF-IDF's sparse keyword-based features provide no discriminative power in this domain. In low-volume, high-noise political text with rapidly evolving vocabulary, sparse features are too brittle: rare and specialized terms dominate the feature space but appear infrequently enough that TF-IDF statistics never stabilize, and the Naive Bayes classifier defaults to majority-class prediction throughout.

\begin{figure*}
    \centering
    \includegraphics[width=\textwidth]{Report/baseline.png} 
    \caption{TF-IDF + MultinomialNB performance on 890 GDELT articles. The model's rolling accuracy (blue solid line) overlaps precisely with the majority-class baseline (red dashed line), indicating that the classifier has collapsed to majority-class prediction without extracting any discriminative signal from article content. ADWIN fails to detect concept drift around the November 5, 2024 U.S. Presidential Election and January 20, 2025 Inauguration (purple and orange vertical lines), a consequence of the model's degenerate behavior.}
    \label{fig:tfidf}
\end{figure*}

\paragraph{LLM Embeddings Word Representation}
Replacing TF-IDF with \texttt{StreamSentenceTransformer} (all-MiniLM-L6-v2) and SGD Logistic Regression produces performance nearly indistinguishable from the majority-class baseline. Results (Figure~\ref{fig:llmreal}) show that the LLM + SGD model (green line) closely tracks the majority-class baseline (red dashed) throughout the stream, with the two curves overlapping extensively across the entire evaluation period. Unlike TF-IDF's stable 0.70--0.75 accuracy, the LLM model exhibits higher variance, but this volatility does not correlate with known political events. 



\begin{figure*}
    \centering
    \includegraphics[width=\textwidth]{Report/llm.png} 
    \caption{LLM Embeddings + SGD pipeline (green) on 890 GDELT articles. The Logistic Regression model closely tracks the rolling majority-class baseline (red dashed) throughout, showing no substantial predictive advantage or detectable response to major political events (November 5, 2024 election and January 20, 2025 inauguration, marked by purple vertical lines). The model's fluctuations appear uncorrelated with these structural changes, suggesting that semantic embeddings alone do not provide sufficient signal for drift detection in this challenging domain.}
    \label{fig:llmreal}
\end{figure*}

\subsection{Unsupervised Distributional Drift Detection}


\paragraph{TF-IDF vs. Embeddings on Mixed Stream}
Daily cosine distances are computed between consecutive window aggregates using both TF-IDF (Figure~\ref{fig:h1_tfidf_mixed}) and \texttt{all-MiniLM-L6-v2} embeddings (Figure~\ref{fig:h1_embeddings_mixed}). TF-IDF produces a cleaner drift score signal with fewer spurious fluctuations, while embeddings exhibit high variance with numerous false-positive alerts. Applying KSWIN drift detection, TF-IDF detects drifts that align with major political events (election, inauguration, trials, conventions), whereas embeddings produce excessive noise without clear event correlation. This suggests TF-IDF's sparse, explicit feature representation is better suited than dense embeddings for detecting distributional shifts in aggregated political text.

\paragraph{Ideology-Specific Drift Patterns}
Analyzing separate left (Figure~\ref{fig:h1_tfidf_left}) and right (Figure~\ref{fig:h1_tfidf_right}) streams reveals asymmetric event sensitivity. The left-leaning stream shows sharp drift signals aligned with: (1) Biden's withdrawal from the race (July 21, 2024), (2) Harris becoming Democratic nominee (August 5, 2024), and (3) Democratic Convention (August 19, 2024). The right-leaning stream exhibits pronounced drifts around: (1) Trump conviction (May 30, 2024), (2) Republican Convention (July 15, 2024), and (3) Trump assassination attempt (July 13, 2024). These patterns suggest that ideologically-specific events produce detectable semantic shifts within their corresponding media streams.

\begin{figure*}
    \centering
    \includegraphics[width=\textwidth]{Report/h1_tfidf_daily_mixed.png} 
    \caption{Part II: TF-IDF daily drift detection on mixed stream. Blue dots: daily cosine distances between consecutive aggregated TF-IDF vectors. Red dashed lines: KSWIN drift alarms. Gray vertical lines: major political events. Mixed stream shows moderate event alignment with concentrated alerts around conventions (Jul-Aug 2024) and election period (Nov 2024).}
    \label{fig:h1_tfidf_mixed}
\end{figure*}

\begin{figure*}
    \centering
    \includegraphics[width=\textwidth]{Report/h1_embedding_daily_mixed.png} 
    \caption{Part II: Dense embedding daily drift detection on mixed stream. Blue dots: daily cosine distances between consecutive aggregated \texttt{all-MiniLM-L6-v2} embeddings. Red dashed lines: KSWIN drift alarms. Gray vertical lines: major political events. Embeddings produce noisier signals with more spurious alerts and weaker event alignment compared to TF-IDF, demonstrating TF-IDF superiority for distributional drift detection in political media.}
    \label{fig:h1_embeddings_mixed}
\end{figure*}

\begin{figure*}
    \centering
    \includegraphics[width=\textwidth]{Report/h1_tfidf_daily_left.png} 
    \caption{Part II: TF-IDF daily drift detection on left-leaning stream. Blue dots: daily cosine distances. Red dashed lines: KSWIN drift alarms. Gray vertical lines: major political events. Left stream exhibits sharp, concentrated drift signals strongly aligned with Democratic-specific events: Biden's withdrawal (July 21, 2024), Harris nomination (August 5, 2024), and Democratic Convention (August 19, 2024), demonstrating ideology-specific event sensitivity.}
    \label{fig:h1_tfidf_left}
\end{figure*}

\begin{figure*}
    \centering
    \includegraphics[width=\textwidth]{Report/h1_tfidf_daily_right.png} 
    \caption{Part II: TF-IDF daily drift detection on right-leaning stream. Blue dots: daily cosine distances. Red dashed lines: KSWIN drift alarms. Gray vertical lines: major political events. Right stream shows pronounced drifts synchronized with Republican-specific events: Trump assassination attempt (July 13, 2024), Republican Convention (July 15, 2024), and Trump conviction (May 30, 2024), demonstrating strong ideology-specific event sensitivity.}
    \label{fig:h1_tfidf_right}
\end{figure*}

%------------------------------------------------
\newpage
\section{Discussion}

\subsection{Key Findings}

\textbf{Part I: Adaptive Streaming
Classifiers.} Both approaches fail catastrophically on real-world political text classification. TF-IDF + MultinomialNB collapses entirely to majority-class prediction, providing zero learning signal. The model's rolling accuracy perfectly overlaps with the trivial baseline, indicating that sparse features are insufficient for this domain: rare specialized political vocabulary dominates the discriminative feature space but appears too infrequently in a 890-article stream for TF-IDF weights to stabilize. LLM embeddings similarly perform at baseline without detectable improvement. Neither approach exhibits correlation with the November 2024 election or January 2025 inauguration.

The synthetic experiments confirm that ADWIN can reliably detect large, abrupt injected drifts under controlled conditions. However, this capability does not transfer to the real-world domain, where both vectorization approaches fail at the more fundamental task of extracting any predictive signal whatsoever. The results highlight a critical challenge: drift detection methods presuppose that a baseline model possesses some discriminative power; when baseline performance collapses to majority class, drift detection becomes irrelevant.

\textbf{Part II: Unsupervised
Distributional Drift Detection.} The unsupervised distributional approach succeeds where classification-based methods fail. On mixed streams, TF-IDF produces superior drift detection compared to dense embeddings, with alerts aligning with some major political events (election, conventions). Critically, analyzing left and right streams separately reveals ideology-specific drift patterns: the left stream detects drifts synchronized with Democratic-relevant events (Biden withdrawal, Harris nomination), while the right stream captures Republican-specific events (Trump conviction, assassination attempt). This demonstrates that distributional drift detection can extract meaningful political signals when analyzed through ideological lenses, even when combined streams show weaker alignment.

\textbf{Unified conclusion:} Part I (supervised classification-based drift) fails catastrophically, with both TF-IDF and LLM approaches collapsing to majority-class prediction. Part II (unsupervised distributional drift) succeeds by abandoning the classification requirement. 

\subsection{Limitations}

The acquired article corpus is too small and restricted to a single left-leaning outlet pair and one right-leaning outlet. The binary left/right categorization is oversimplified; real media bias varies by outlet, author, and topic. Part I's online classifiers may not be suited for the low signal and high noise of political text; Part II's drift detectors may be miscalibrated to the noise level of cosine distance streams. LLM embeddings use a frozen pre-trained model not fine-tuned on political discourse.

\subsection{Future Work}

Given the results from both supervised and unsupervised approaches, future work could investigate: (1) whether larger article volumes (multiple outlets per ideology, longer time span) reveal clearer drift signals; (2) alternative vectorization approaches (e.g., fine-tuned domain-specific embeddings, topic modeling, lexicon-based sentiment/framing metrics);  A supervised fine-tuned language model (e.g., RoBERTa on political discourse) might improve Part I classification, enabling more reliable drift detection. 

%----------------------------------------------------------------------------------------
%	 REFERENCES
%----------------------------------------------------------------------------------------

\printbibliography

%----------------------------------------------------------------------------------------

\end{document}